\documentclass[11pt]{article}
\usepackage{amsmath, amssymb}
\usepackage{geometry}
\geometry{margin=0.7in}
\setlength{\parindent}{0pt}
\setlength{\parskip}{0.35em}

\begin{document}
\begin{center}
{\LARGE Rational Functions Cheatsheet}\\[0.4em]
{\large Fundamental Review}
\end{center}

\section*{Definition}
\[
f(x)=\frac{p(x)}{q(x)},\qquad q(x)\ne 0
\]

\section*{Domain}
Exclude values that make the denominator zero.

\section*{Holes vs Vertical Asymptotes}
\begin{itemize}
  \item Common factor cancels $\Rightarrow$ hole
  \item Denominator zero that does not cancel $\Rightarrow$ vertical asymptote
\end{itemize}

\section*{Intercepts}
\[
y\text{-intercept}: f(0)
\]
\[
x\text{-intercepts}: p(x)=0 \text{ after cancellation, with denominator nonzero}
\]

\section*{Horizontal Asymptotes}
For
\[
f(x)=\frac{a_nx^n+\cdots}{b_mx^m+\cdots}
\]
\begin{itemize}
  \item if $n<m$, then $y=0$
  \item if $n=m$, then $y=\dfrac{a_n}{b_m}$
  \item if $n>m$, no horizontal asymptote
\end{itemize}

\section*{Slant Asymptote}
If degree of numerator is exactly one more than degree of denominator, divide:
\[
\frac{p(x)}{q(x)}=\text{quotient}+\frac{\text{remainder}}{q(x)}
\]
The quotient line is the slant asymptote.

\section*{Sign Chart Idea}
Use zeros and vertical asymptotes as boundary points. Test the sign of each interval.

\section*{Quick Graph Checklist}
\begin{itemize}
  \item factor numerator and denominator
  \item identify domain restrictions
  \item check for holes
  \item find vertical and horizontal/slant asymptotes
  \item find intercepts
  \item use sign chart behavior between key $x$-values
\end{itemize}

\end{document}
